\section{A single Hecke module for all the split columns}

% BEGIN G82C1
\paragraph{G82C1. A simultaneous presentation with finite local cocycles.}
Fix the guarded split middle-range index and good point of G81S1.
Work over the common good coefficient extension of G81S3, with fraction
field $L$, DVR $O$ and $k_O=\operatorname{length}_O(O/p)$.
Write $X'\to X_i$ for the ordered-root cover, $A$ for its integral
matrix order, and $T'$ for its central Hecke scalars. Let $E$ be the
first-character matrix and $e$ the rank-one idempotent in the rank
corner. Thus $E^2=\delta E$ with $\delta\in T'$ nonzero at every
point, and $e^2=e$. At the specified lift $\pi'$ put
$\delta_0=\delta(\pi')$, $v=v_p(\delta_0)=v_\delta(i)$, and
$a=X.e_S(i)$. Put $U_g=\chi_1(g)^{-1}R(g)$, $y_g=(U_g-1)E$,
and $Y=\sum_g A y_g$. The multiplication in these formulas is in
the same matrix order at all places.

For a column $u$ write $D_u$ for its decomposition group,
$T$ for the chosen rank lattice in $W$, and
$j_u=\dim_L H_f^1(D_u,W)$. Take integral cocycles
$\kappa_{uh}$, $1\le h\le j_u$, whose rational classes are an
$O$-basis of the image of $H_f^1(D_u,T)$ in $H^1(D_u,W)$.
Fix a generator $v_1$ of the direct summand $eT$. F82L1--F82L4
supply finite cocycles $K_{uh}:D_u\to Ae$ and $k_{u0}:D_u\to AE$
for the twisted left action. On $v_1$ at the rank block the
first cocycle has class $p^a[\kappa_{uh}]$. We may impose equality
of cocycles there: subtract the coboundary of any lift under the
surjection $Ae\to T$ of the integral correcting vector. At the
first character the second cocycle has value
$p^a\delta_0\,\deg_u(g)$, where the continuous unramified degree
has value $1$ at a Frobenius lift. That line has trivial twisted
action and hence no coboundary ambiguity.

Form the left $A$-module $N$ by adjoining to $Y$, for each $u$,
the generators $q_u,s_{u1},\ldots,s_{u j_u},b_u$. Impose
\[
 E q_u=0,\qquad (1-e)s_{uh}=0,\qquad (\delta-E)b_u=0,
\]
and, for every $g\in D_u$, impose
\[
 y_g+(U_g-1)q_u
    =\delta\left(\sum_{h=1}^{j_u}K_{uh}(g)s_{uh}
                          +k_{u0}(g)b_u\right).
\]
These are precisely the local relations of $\mathrm{ColNNJ}$.
For any left $A$-module $M$ a map $N\to M$ amounts to an
$A$-map $\varphi:Y\to M$ and images $Q_u,S_{uh},B_u$ satisfying
these same equations. This follows by first taking the direct sum
of $Y$ with the free modules on the indicated generators, then
quotienting by the left submodule spanned by the relations.
In particular there is only one $\varphi$ for all columns. The order
$A$ is finite over $O$ inside the product of integral matrix rings.
Thus $Y$ is finite over $O$ and $N$ admits a finite presentation over
$A$: a finite free module over $A$ surjects onto it, and the kernel
is an $O$-submodule of a finite $O$-module.
% END G82C1

% BEGIN G82C2
\paragraph{G82C2. Nonzero evaluations and the determinant index.}
For $x'$ above $x\notin Z_i$ let $V=V_{x'}$ have the twisted
action $U_g$, with dimension $n_V=r+t$. We construct a nonzero
$A$-map from $N$ to $V$. Choose a parameter $\eta\in E(x')V$
and take $\varphi(y)=y(x')\eta$. Its values on the generators
are $\delta(x')(U_g-1)\eta$. Since $V$ has no global fixed vector
by D79G7 and $\delta(x')\ne0$, $\eta\mapsto\varphi$ is injective.
Independently at each column choose
\[
 Q_u\in\ker E(x'),\qquad S_{uh}\in e(x')V,\qquad
 B_u\in E(x')V=\ker(\delta(x')-E(x')).
\]
The three generator relations hold. Both image lines have dimension
one. Together with $\eta$ the parameter space therefore has
dimension $1+\sum_u(n_V+j_u)$.

Send the parameters to the left minus the right side of the last
relation of G82C1, as functions on the several $D_u$. At $u$ the
image lies in the space of continuous rational finite cocycles
$Z_f^1(D_u,V)$. The two terms on the left are coboundaries, and
on the right evaluation on the selected vectors is an equivariant
map out of the rationalized left ideal, so preserves finiteness by
F82L1--F82L4. At this column the same results give
$\dim H_f^1(D_u,V)=j_u+\dim V^{D_u}$.
Coboundaries have dimension $n_V-\dim V^{D_u}$, so the target
space has dimension $n_V+j_u$. Rank-nullity for the simultaneous
map gives a nonzero kernel vector. Its tuple map is nonzero since
the generator values recover all local parameters and the map from
$\eta$ is injective. This proves the required assertion at every
point over the complement using the same module $N$.

For a distinguished block $b$ of dimension $d_b$, write
$M_b=(L/O)^{d_b}$ with its block action, and $H_b=\operatorname{Hom}_A(N,M_b)$.
Then for the original evaluation ideal $I_\pi$ extended to $O$,
\[
                  \operatorname{length}_O(O/I_\pi)
                          \le \sum_b \operatorname{length}_O H_b .
\]
To see this, present $N$ on $s$ generators. For any square matrix
of relations on them, the $s$ images of these generators under a
nonzero map just constructed give a nonzero kernel vector of the
evaluated matrix off support. This holds simultaneously for all
relation matrices. G79E11's determinant argument applies: for an
indeterminate linear combination of any finite list of such matrices
all coefficients of the block determinant are Hecke scalars, and
vanishing on every ordered lift over the complement puts their
values at $\pi'$ in $I_\pi$. The free ordered-root basis gives
this descent, including for the perturbed corner of G81S3.

For clarity, evaluate all relation rows in block $b$ and take the
span $P_b$ of their individual matrix rows inside $O^{sd_b}$.
The solution space is exactly
$H_b=\operatorname{Hom}_O(O^{sd_b}/P_b,L/O)$ by the presentation.
If any $P_b$ has deficient rank the displayed inequality follows
with infinite right side. Otherwise Smith reduction identifies its
basis determinant valuation with $\operatorname{length}_O H_b$.
Surjectivity $A\to M_{d_b}(O)$ lets one put each basis row in any
chosen row position by multiplying relations by lifts of matrix
units and adding. This makes one square relation matrix for each
block with the specified basis determinant there. In independent
variables take their linear combination and factor its determinant
at $\pi'$ into the block determinants. In block $b$ its minimum
coefficient valuation is at most the basis determinant valuation,
by specializing the variables to select its matrix. Those minima
add in a product since nonzero residue polynomials have nonzero
product. Every coefficient of the product belongs to $I_\pi$, giving
the inequality as in G79E12. This polynomial argument also works
over a small residue field.
% END G82C2

% BEGIN G82C3
\paragraph{G82C3. The rank block is finite at all columns.}
For a tuple in $H_b$ put $z(g)=\varphi(y_g)$. The identity
$y_{gg'}=y_g+U_g y_{g'}$ gives its cocycle law.
The map is continuous: $Y$ is a finite $O$-lattice and its
subspace topology from the matrix family agrees with its lattice
topology (Smith form over the DVR); $\varphi$ has bounded torsion
image and open kernel. The restriction of $\varphi$ to the
generators of $Y$ is determined by $z$. Moreover $z$ lies in the
block cocycle module used in G81S3--G81S4: extend $\varphi$
$O$-linearly to the integral matrix family by divisibility of
$M_b$, retaining equivariance on $Y$ as in G79E12--G79E13.
Consequently the normalized loss in mapping the image of $z$ to
global cohomology is at most $2d_b v$ for every block except the
first character, where it is zero (G81S3). Here and below lengths
are divided by $k_O$.

On the rank block $M=W/T$ the matrix $E$ acts by zero and $eM$
is the line generated over $L/O$ by $v_1$, meaning the image of
$t v_1$ depends only on $t\bmod O$. Write
$S_{uh}=s_{uh}v_1$, $s_{uh}\in L/O$. We have
$B_u\in M[\delta_0]$ and $k_{u0}(g)$ acts by zero because
it lies in $AE$. The local relation now reads
\[
 z(g)+(U_g-1)Q_u
       =p^a\delta_0\sum_h \kappa_{uh}(g)s_{uh}\quad(g\in D_u)
\]
in $M$, with the product on the right computed using any lifts of
the scalars to $L$. Its rational lift is a finite cocycle, so the
local class of $z$ lies in $H_f^1(D_u,M)$ at every column.

Consider the kernel of the map to actual cocycles $z$. When $z=0$,
the right side has zero class in $H^1(D_u,M)$. Its rational
finite class is then in the image of $H^1(D_u,T)$, by exactness
for $0\to T\to W\to M\to0$. The integral preimages are finite
by the definition of the integral finite part. The basis chosen
in G82C1 therefore forces $p^a\delta_0 s_{uh}=0$ in $L/O$
for every $h$. These parameters have total normalized length at
most $j_u(a+v)$. The choices for $B_u$ have length $r v$.
For a fixed choice of these parameters, any two possible $Q_u$
differ by an invariant of $M$, of normalized length
$h_W(i,u,\pi)$. Equivalently the kernel of projection to the
$s_{uh}$ and $B_u$ embeds in the product of these invariant
modules over $u$. Thus the entire normalized kernel has length
at most
\[
             \sum_u \{j_u(a+v)+rv+h_W(i,u,\pi)\}.
\]
The estimate includes $j_u=0$.
% END G82C3

% BEGIN G82C4
\paragraph{G82C4. The character blocks and the unramified degree.}
In character block $j$ with $0<j<t$, $M=(L/O)(-j)$ and
$E=e=0$. Therefore $S_{uh}=0$, $B_u\in M[\delta_0]$, and
$k_{u0}(g)$ acts by zero. The relation is
$z(g)=-(U_g-1)Q_u$ on $D_u$. Its local class is zero at every
column. If $z=0$ then $Q_u\in M^{D_u}$. The kernel of the map
to actual cocycles thus embeds in the product of these invariants
and the $\delta_0$-torsion lines. Its normalized length is at
most $\sum_u(h_0+v)$, where $h_0$ is the maximum normalized
invariant length of the nontrivial character ratios at the columns.

On the first character $M=L/O$ the action $U_g$ is trivial,
$e=0$, and $E=\delta_0$. Here $Q_u\in M[\delta_0]$,
$S_{uh}=0$, and the relation on $D_u$ becomes
\[
                z(g)=p^a\delta_0^2\deg_u(g) B_u.
\]
This gives an unramified class at every column. In the kernel
$z=0$ a Frobenius lift with degree $1$ shows
$p^a\delta_0^2 B_u=0$. The generator values now embed the kernel
in a product of scalar torsion modules of normalized lengths $v$
and $a+2v$ at each column. The resulting bound is
$\sum_u(a+3v)$. Both occurrences of $\delta_0$ in the last
formula are needed: one comes from the defining relation and the
other from evaluating the lifted cocycle in $AE$.
% END G82C4

% BEGIN G82C5
\paragraph{G82C5. The guarded split middle-range inequality.}
Let $c$ be the number of split columns. Add G82C3--G82C4 over
the $r$-dimensional rank block and the $t$ individual character
blocks. Their total normalized loss from tuple maps to cohomology
classes is at most
\[
 \begin{aligned}
 &\sum_u\{a+j_u(a+v)+h_W(i,u,\pi)+(n_V+2)v\}\\
 &\hspace{12mm} +(t-1)c h_0+2(n_V-1)v .
 \end{aligned}
\]
Indeed the $v$ terms per column outside $j_u(a+v)$ have
coefficients $r,t-1,3$ respectively, and the last term is the
global coboundary bound of G81S3. Since $j_u\le d_u m$ by G81S1
and $n_V c+2(n_V-1)\le4n_V(c+1)$, the loss is bounded by
\[
 \sum_u\{a+d_um(a+v)+2v+h_W(i,u,\pi)\}
                 +tc h_0+4n_V(c+1)v.
\]

All the resulting rank classes are finite at the columns. At each
conjugate row $w$ their image modulo the finite part has normalized
length at most $d_um a+rX.c_0(i,u)$ by the graph construction and
Tate pairing of I81G3--I81G4 and the row bound G81S2. These results
apply to any finite list of the present cocycles by their cut
presentation in G82C3. A uniform length bound on the spans of
finite lists bounds the full span by I81G2. G81S4 supplies all
other local conditions for the remaining rank classes to embed
in the native Selmer group, of normalized length $s_\pi$.

Each character has its column condition by G82C4. At each row the
loss in imposing its finite part is at most $(d_w+1)X.c_B(i)$
by M81B3 and I81G5, and at the auxiliary inertias it is at most
$\sum_{P\in Q_i}z_P/k_O$ by G79E8. The remaining trivial-character
classes are everywhere unramified and the remaining nontrivial
character classes vanish locally above $p$. G79E6 and G81S4 put
their normalized lengths at most $v_p(h_E)$ and the native $u_j$
for $0<j<t$, respectively. These restrictions on submodules of
the global block classes use the same away-from-$p$ hypotheses as
G81S4. In particular summing the auxiliary contributions over the
$t$ characters costs at most the killing allowance $p^{-2}$.

Insert all these bounds in G82C2. Enlarge $d_w+1$ to $d_w+2$
in the row character allowance. For each conjugate pair the
column and row terms in $a$ now give $(2d_um+1)a$, with remaining
terms $d_um v+2v+h_W(i,u,\pi)+rX.c_0(i,u)$. The other errors
are $tX.c_B(i)\sum_w(d_w+2)$, $tc h_0$ and $4n_V(c+1)v$.
By the split formulas of DFD8 and N78T2 these are exactly the
stored $\mathrm{len0}$ and $\mathrm{errorBoundHub}$ allowances
as grouped in G81S4. Normalized lengths descend to the original
coefficient DVR by I81G2 and G81S3: both finite lengths and
$k_O$ scale by the relative ramification index, including the
evaluation ideal, and the Selmer and local cohomology conditions
commute with this finite free extension. Under the original good
point, killing and odd finiteness guards this proves the split
middle-range $\mathrm{StarAD}$ inequality of the chosen hub.
The determinant applied here is the one for $N$; the larger
restriction image introduced in C81S2 is unnecessary for this
inequality.
% END G82C5
